%% guide-bouilloire-courbe — chauffage de 1,00 kg d'eau dans une bouilloire électrique.
%% Portion linéaire de (0 s ; 20 °C) à (250 s ; 100 °C) puis palier d'ébullition à 100 °C.
%% Échelles : 1 s -> 0,0186 cm ; 1 °C -> 0,0283 cm.
%% RÈGLE : la pente (coefficient directeur) n'est pas écrite — c'est la réponse de l'étape 3.
\documentclass[tikz,border=4pt]{standalone}
\usepackage{liaisons}
\begin{document}
\begin{tikzpicture}[x=1cm,y=1cm,
  axe/.style={-{Stealth[length=5pt]}, line width=0.6pt},
  courbe/.style={draw=solideB, line width=1.4pt, line join=round, line cap=round},
  grille/.style={line width=0.3pt, black!15, dotted},
  rappel/.style={line width=0.5pt, black!50, dash pattern=on 2.5pt off 2pt},
  grad/.style={font=\scriptsize, inner sep=2.5pt},
  tri/.style={{Stealth[length=4.5pt]}-{Stealth[length=4.5pt]}, line width=0.9pt, draw=solideA}]

\newcommand\xt[1]{{(#1)*0.0186}}      % abscisse d'une date en s
\newcommand\yt[1]{{(#1)*0.0283}}      % ordonnée d'une température en °C

%% ================= quadrillage ============================================
\foreach \t in {50,100,...,350} { \draw[grille] (\xt{\t},0) -- (\xt{\t},\yt{120}); }
\foreach \q in {20,40,...,120}  { \draw[grille] (0,\yt{\q}) -- (\xt{350},\yt{\q}); }

%% ================= axes et graduations ====================================
\draw[axe] (-0.12,0) -- (\xt{378},0);
\node[grad, anchor=south east] at (\xt{380},0.04) {$t$ (s)};
\draw[axe] (0,-0.12) -- (0,\yt{132});
\node[grad, anchor=south west] at (0.04,\yt{132}) {$\theta$ ($^\circ$C)};
\node[grad, anchor=north east] at (-0.02,-0.02) {$O$};
\foreach \t in {50,100,150,200,250,300,350} {
  \draw[line width=0.5pt] (\xt{\t},0) -- (\xt{\t},-0.09);
  \node[grad, anchor=north] at (\xt{\t},-0.09) {\t}; }
\foreach \q in {20,40,60,80,100,120} {
  \draw[line width=0.5pt] (0,\yt{\q}) -- (-0.09,\yt{\q});
  \node[grad, anchor=east] at (-0.09,\yt{\q}) {\q}; }

%% ================= traits de rappel =======================================
\draw[rappel] (0,\yt{100}) -- (\xt{320},\yt{100});
\draw[rappel] (\xt{250},0) -- (\xt{250},\yt{100});

%% ================= le tracé ===============================================
\draw[courbe] (\xt{0},\yt{20}) -- (\xt{250},\yt{100}) -- (\xt{320},\yt{100});
\node[font=\scriptsize, text=solideB, anchor=south] at (\xt{287},\yt{103})
  {palier d'ébullition};

%% ================= triangle du coefficient directeur ======================
\draw[tri] (\xt{6},\yt{20}) -- (\xt{244},\yt{20});
\node[font=\small, text=solideA, fill=white, inner sep=1.5pt]
  at (\xt{125},\yt{20}) {$\Delta t$};
\draw[tri] (\xt{250},\yt{24}) -- (\xt{250},\yt{96});
\node[font=\small, text=solideA, fill=white, inner sep=1.5pt]
  at (\xt{250},\yt{60}) {$\Delta\theta$};

%% ================= cartouche des données ==================================
\node[draw=black!55, line width=0.6pt, rounded corners=2pt, fill=black!3,
      inner sep=4pt, align=center, font=\scriptsize, anchor=north west]
  at (\xt{75},\yt{158})
  {$m = 1{,}00$ kg \quad $R = 40\ \Omega$ \quad $U_{\text{eff}} = 230$ V\\[1pt]
   bouilloire très bien isolée};
\end{tikzpicture}
\end{document}
